\documentclass[11pt,a4paper]{article}
\usepackage{fontspec}
\usepackage{kotex}
\usepackage{geometry}
\usepackage{graphicx}
\usepackage{xcolor}
\usepackage{tcolorbox}
\usepackage{booktabs}
\usepackage{array}
\usepackage{longtable}
\usepackage{colortbl}
\usepackage{fancyhdr}
\usepackage{titlesec}
\usepackage{hyperref}
\usepackage{emoji}
\usepackage{amsmath}
\usepackage{amssymb}
\usepackage{enumitem}

\geometry{left=25mm, right=25mm, top=30mm, bottom=30mm}

\setmainfont{Noto Sans CJK KR}
\setsansfont{Noto Sans CJK KR}
\setmonofont{Noto Sans CJK KR}
\setemojifont{Noto Color Emoji}

\definecolor{primary}{HTML}{1976D2}
\definecolor{darkbrown}{HTML}{3E2723}
\definecolor{mediumbrown}{HTML}{5D4037}
\definecolor{lightbrown}{HTML}{8D6E63}
\definecolor{cream}{HTML}{FFF3E0}
\definecolor{lightgray}{HTML}{F5F5F5}
\definecolor{tableheader}{HTML}{E8E8E8}

\titleformat{\section}
  {\Large\bfseries\color{primary}}
  {\thesection.}{0.5em}{}
\titleformat{\subsection}
  {\large\bfseries\color{darkbrown}}
  {\thesubsection}{0.5em}{}
\titleformat{\subsubsection}
  {\normalsize\bfseries\color{mediumbrown}}
  {\thesubsubsection}{0.5em}{}

\renewcommand{\contentsname}{목차}
\renewcommand{\figurename}{그림}
\renewcommand{\tablename}{표}

\tcbset{
  tipbox/.style={colback=cream,colframe=primary,boxrule=1.5pt,arc=3mm,
    left=5mm,right=5mm,top=3mm,bottom=3mm,fonttitle=\bfseries\color{primary}},
  formulabox/.style={colback=lightgray,colframe=lightgray,boxrule=0pt,arc=2mm,
    left=5mm,right=5mm,top=3mm,bottom=3mm},
  summarybox/.style={colback=white,colframe=primary,boxrule=2pt,arc=4mm,
    left=5mm,right=5mm,top=5mm,bottom=5mm}
}

\hypersetup{colorlinks=true,linkcolor=primary,urlcolor=primary,bookmarks=true,unicode=true}


\pagestyle{fancy}
\fancyhf{}
\fancyhead[R]{\small\color{lightbrown}마이크로비트 센서 검색 도구 이론해설서 v1.0.0.0}
\fancyfoot[C]{\thepage}
\renewcommand{\headrulewidth}{0.4pt}

\begin{document}

\begin{titlepage}
\centering
\vspace*{3cm}
{\Huge\bfseries\color{primary}\emoji{magnifying-glass-tilted-left} 마이크로비트 센서 검색 도구}\\[0.5cm]
{\Huge\bfseries\color{primary} 이론해설서}
\\[1.5cm]
{\Large 시뮬레이터 버전: v1.0.3.0}
\\[2cm]
\begin{tcolorbox}[summarybox]
\centering
{\large 이 이론해설서는 마이크로비트 센서 검색 도구이 다루는 핵심 개념과 원리를 설명합니다.}
\end{tcolorbox}
\vfill
{\large 사물인터넷과 빅데이터}\\[0.3cm]
{\large 청주교육대학교 컴퓨터교육과}
\vspace{1cm}
\end{titlepage}

\section*{1. 이 문서의 목적}

이 이론해설서는 \textbf{마이크로비트 센서 검색 도구}이 다루는 핵심 개념과 원리를 설명합니다. 시뮬레이터를 사용하기 전에 배경 지식을 쌓거나, 사용 후 깊이 있는 이해를 위해 참조합니다.


\begin{tcolorbox}[summarybox]
\textbf{핵심 주제}: 센서 선택과 프로젝트 설계

\end{tcolorbox}


\section*{2. 핵심 개념}

\begin{center}
\renewcommand{\arraystretch}{1.4}
\begin{tabular}{|p{3.1cm}|p{11.9cm}|}
\hline
\rowcolor{tableheader}
\textbf{개념} & \textbf{설명} \\
\hline
\textbf{측정 가능성} & 센서로 직접 측정할 수 있는 물리량. '집중도'는 직접 측정 불가, '온도'는 가능 \\
\hline
\textbf{프록시 변수} & 직접 측정할 수 없는 것을 대신 나타내는 간접 측정값 (예: 집중도 \textrightarrow{} 자세·시선·소음) \\
\hline
\textbf{센서 조합 최적화} & 목표 데이터를 수집하기 위한 최소한의 센서 조합을 찾는 과정 \\
\hline
\textbf{내장 vs 외장 센서} & 마이크로비트에 내장된 기본 센서 vs Grove로 확장하는 정밀 센서 \\
\hline
\textbf{연구 질문} & 프로젝트의 출발점이 되는 핵심 질문. '측정 가능한 질문'으로 변환 필요 \\
\hline
\end{tabular}
\end{center}


\section*{3. 원리 상세}

\subsection*{측정 가능한 연구 질문 만들기}

IoT 프로젝트의 첫 단계는 "무엇을 측정할 것인가?"를 결정하는 것입니다.


\textbf{좋은 연구 질문의 조건:}

1. \textbf{측정 가능} — 센서로 데이터를 수집할 수 있어야 함

2. \textbf{구체적} — "환경"보다 "교실의 온도와 소음"

3. \textbf{비교 가능} — "높다/낮다"를 판단할 기준이 있어야 함


\textbf{변환 예시:}


\begin{center}
\renewcommand{\arraystretch}{1.4}
\begin{tabular}{|p{4.4cm}|p{10.6cm}|}
\hline
\rowcolor{tableheader}
\textbf{\emoji{cross-mark} 측정 불가능} & \textbf{\emoji{check-mark-button} 측정 가능} \\
\hline
"우리 교실은 쾌적한가?" & "교실 온도·습도·CO\textsubscript{2}가 쾌적 기준 이내인가?" \\
\hline
"학생들이 집중하는가?" & "소음 수준과 교실 온도가 집중에 적합한 범위인가?" \\
\hline
"날씨가 좋은가?" & "기압·온도·습도 변화 추이가 맑은 날 패턴과 일치하는가?" \\
\hline
\end{tabular}
\end{center}

\subsection*{센서 선택 의사결정}

센서 선택 시 고려할 요소:

1. \textbf{필수 데이터}: 연구 질문에 반드시 필요한 측정값

2. \textbf{내장 vs 외장}: 내장 센서로 충분하면 추가 비용 없음

3. \textbf{정밀도}: 대략적 추세면 내장, 정확한 수치면 외장

4. \textbf{비용}: Grove 센서당 5,000\textasciitilde{}20,000원



\section*{4. 실생활 응용 사례}

\textbf{사례 1}: 스마트 빌딩: 건물 관리에 필요한 센서 종류·위치 계획


\textbf{사례 2}: 환경 모니터링: 하천 수질 측정을 위한 센서 조합 설계


\textbf{사례 3}: 농업 IoT: 작물별 필요 환경 데이터에 맞는 센서 배치


\textbf{사례 4}: 산업 IoT: 공장 설비 이상 감지를 위한 센서 선정



\section*{5. 더 알아보기}

$\bullet$ \textbf{센서 데이터시트}: 제조사가 제공하는 센서 사양서 읽는 법

$\bullet$ \textbf{센서 퓨전 알고리즘}: 여러 센서 데이터를 결합하는 고급 기법

$\bullet$ \textbf{연구 설계}: 가설 \textrightarrow{} 변수 정의 \textrightarrow{} 데이터 수집 \textrightarrow{} 분석의 연구 방법론

$\bullet$ \textbf{측정 불확실도}: 모든 측정에는 오차가 있다는 개념



\section*{6. 핵심 용어 사전}

\begin{center}
\renewcommand{\arraystretch}{1.4}
\begin{tabular}{|p{2.6cm}|p{2.6cm}|p{9.8cm}|}
\hline
\rowcolor{tableheader}
\textbf{용어} & \textbf{학술 용어} & \textbf{쉬운 설명} \\
\hline
측정 가능성 & — & 센서로 직접 측정할 수 있는 물리량 \\
\hline
프록시 변수 & — & 직접 측정할 수 없는 것을 대신 나타내는 간접 측정값 (예: 집중도 \textrightarrow{} \\
\hline
센서 조합 최적화 & — & 목표 데이터를 수집하기 위한 최소한의 센서 조합을 찾는 과정 \\
\hline
내장 vs 외장 센서 & 내장 vs 외장 센서 & 마이크로비트에 내장된 기본 센서 vs Grove로 확장하는 정밀 센서 \\
\hline
연구 질문 & — & 프로젝트의 출발점이 되는 핵심 질문 \\
\hline
\end{tabular}
\end{center}


\section*{7. 참고자료}

\begin{center}
\renewcommand{\arraystretch}{1.4}
\begin{tabular}{|p{1.5cm}|p{7.1cm}|p{6.4cm}|}
\hline
\rowcolor{tableheader}
\textbf{\#} & \textbf{자료명} & \textbf{비고} \\
\hline
1 & 마이크로비트센서검색도구\_퀵가이드 & 소프트웨어 사용 중 빠른 참조 \\
\hline
2 & 마이크로비트센서검색도구\_사용설명서 & 전체 기능 상세 \\
\hline
3 & 마이크로비트센서검색도구\_활용가이드 & 수업 적용 \\
\hline
4 & 마이크로비트센서검색도구\_개발참조 & 기술 사양 \\
\hline
\end{tabular}
\end{center}


\vfill
\begin{center}
\small\color{lightbrown}
\textit{마이크로비트 센서 검색 도구 이론해설서}
\end{center}
\end{document}